% !TeX program = xelatex
% 2026-09-23：所有修订标红文字与公式统一恢复为正常黑色。
% 历史修订基准保存于 tmp/revision_20260922/before/。
% 研究内容统一为概述与三项子内容；突出机理、方法与验证之间的对应关系。
% 研究目标及两项科学问题均采用单段论述；反向建模术语全文统一。
% 研究方案保留三组核心公式；可行性分析参照指定面上及重大项目材料重写。
% 可编译主文件；参考文献在本文件内，技术路线图采用正文内的TikZ矢量图。
% 正文按申请书模板的一、二、三部分组织。题目另填模板“项目名称”。
% 申请金额按重点课题上限30万元测算；技术指标为拟定目标。
% 已有工作依据：26-0912.pdf、已有组会材料及已核查的控制程序。
% 匿名研究稿不列作已发表论文；机器人代码仅作只读核查。
% 本项目不采用结构模态建模；文献、已有成果与奖励中的相关名称保留原有表述。
% 地面实验依托现有基座振动模拟与机器人作业实验平台，以基座运动模拟柔性桁架振动。
\documentclass[UTF8,a4paper,zihao=-4,fontset=fandol]{ctexart}
\usepackage[left=2.6cm,right=2.6cm,top=2.5cm,bottom=2.5cm]{geometry}
\usepackage{amsmath,amssymb,bm,booktabs,longtable,array,tabularx}
\usepackage{enumitem,setspace,xcolor,caption,needspace,placeins}
\usepackage{graphicx}
\usepackage{tikz}
\usetikzlibrary{arrows.meta,positioning}
\graphicspath{{./}{../}}
\usepackage[unicode,hidelinks]{hyperref}
\usepackage{xurl}
\hypersetup{pdftitle={柔性桁架支撑下空间机器人的运动控制与接触引导组装研究},pdfsubject={南航重点实验室开放课题申请书报告正文}}
\setstretch{1.32}
\setlength{\parindent}{2em}
\setlength{\parskip}{0.15em}
\setlength{\emergencystretch}{2em}
\clubpenalty=10000
\widowpenalty=10000
\displaywidowpenalty=10000
\raggedbottom
\setlist[enumerate]{leftmargin=2.2em,itemsep=0.3em,topsep=0.3em}
\captionsetup{font=small,labelfont=bf,labelsep=quad,skip=8pt}
\ctexset{
 section={format=\zihao{4}\bfseries,name={,、},number=\chinese{section},aftername={},beforeskip=1.2em,afterskip=0.7em},
 subsection={format=\zihao{-4}\bfseries,number=\arabic{subsection},aftername={.\ },beforeskip=1em,afterskip=0.5em},
 subsubsection={format=\zihao{-4}\bfseries,number=\arabic{subsection}.\arabic{subsubsection},aftername={\quad},beforeskip=0.8em,afterskip=0.3em}
}
\newcommand{\topic}[1]{\par\Needspace{4\baselineskip}\addvspace{0.5em}\noindent\textbf{#1}\par\nobreak}
\AddToHook{cmd/section/before}{\Needspace{9\baselineskip}}
\AddToHook{cmd/subsection/before}{\Needspace{7\baselineskip}}
\AddToHook{cmd/subsubsection/before}{\Needspace{5\baselineskip}}
\newcommand{\needinfo}[1]{\par\noindent{\color{black}\textbf{〔待补充〕}#1}\par}
\newcommand{\doi}[1]{\href{https://doi.org/#1}{\nolinkurl{#1}}}
\newcommand{\SE}{\mathrm{SE}(3)}
\newcolumntype{Y}{>{\raggedright\arraybackslash}X}
\newcolumntype{P}[1]{>{\raggedright\arraybackslash}p{#1}}
\renewcommand{\arraystretch}{1.25}
\pagestyle{plain}

\begin{document}

\begin{center}
 {\zihao{3}\bfseries 柔性桁架支撑下空间机器人的运动控制与接触引导组装研究\par
\par}
 \vspace{0.55em}
 {\zihao{-4}航空航天结构力学及控制全国重点实验室开放课题\quad 重点课题}
\end{center}
\vspace{0.35em}

% ==================== 报告正文开始 ====================
\section{立项依据与研究内容}

\subsection{项目的立项依据}

\subsubsection{研究意义}

大型空间天线、太空望远镜等装备的发展，对空间设施的尺度、功能与建造精度提出了更高要求。受运载包络和发射能力限制，依靠单次整体发射与在轨展开难以满足建造需求，通过模块分批运输、机器人在轨转运与装配逐步形成大型设施，成为拓展空间装备能力的重要途径\cite{hu2025,wu2024}。在这一过程中，机器人需要依附已建结构完成模块拾取、转运、接口对准与安装等操作，基座运动、负载变化和接触状态转换使其面临复杂的动力学条件，传统基于固定基座和理想动力学模型的控制方法难以兼顾运动精度与接触柔顺性。因此，亟需发展适应基座运动与多源不确定性的机器人运动与组装控制方法，满足大型空间设施自主、精密、可靠建造的需求。

机器人依附柔性桁架作业时，机械臂运动产生的反作用力与力矩经基座激发桁架振动，桁架振动又引起基座位姿变化并传递至作业端，形成双向动力学耦合；惯性参数、负载及执行时延等偏差进一步影响控制精度。提高反馈增益虽可减小部分运动误差，却限制了机器人顺应接触约束的调整能力，接口位姿偏差与基座振动易引起接触冲击或卡滞。同时，高维耦合模型的在线求解难以兼顾描述精度与反馈实时性。因此，需要揭示基座运动、动力学失配与接触作用对末端响应的影响规律，协调振动扰动下的运动精度与接触柔顺性，实现振动环境下的精确对准与平稳插接。

此外，上述耦合系统具有高维数、多时间尺度和模型不确定性等特征，基于高维模型在线求解结构力学响应，难以兼顾动力学描述精度与反馈控制的实时性。面向在轨作业的测控条件，需要提取影响末端运动与接触响应的主要耦合信息，降低控制对完整结构模型和高维动力学求解的依赖，并结合运动与接触反馈调节控制作用，以适应基座运动及接触状态变化。

针对上述问题，本项目以柔性桁架支撑的空间机器人为研究对象，研究面向实时控制的耦合动力学表征方法，结合几何阻抗设计与动力学优化开展机器人运动控制，利用接触几何信息和接触反馈实现接口对准与柔顺组装，并通过地面试验考核验证。研究有助于阐明柔性桁架振动与多源不确定性共同作用下的末端运动和接触响应规律，形成兼顾实时性、运动精度与接触柔顺性的控制方法，为大型空间天线、太空望远镜等装备的机器人在轨组装提供理论依据与技术支撑。

\subsubsection{国内外研究现状及发展动态}

\topic{（1）柔性结构与空间机器人的耦合动力学}

在轨操作由航天器捕获向大型结构组装拓展，推动动力学研究由基座与机械臂的运动耦合，转向机器人运动、结构变形与组装过程的共同作用\cite{hu2025,wu2024,papadopoulos2021}。机器人所依附结构的质量分布、连接关系和边界随组装进程变化，模型不仅需要描述给定构型下的运动响应，还需要反映构型演化及作业状态转换带来的动力学变化。

现有方法主要在动态特征保留与模型规模之间寻求平衡。模态缩减和模型降阶通过保留主要变形特征降低系统维数\cite{sonneville2021,min2026}；时变边界研究进一步表明，模型参数与约束需随作业状态调整\cite{ni2024}。超大型结构的组装分析揭示了科里奥利效应等空间扰动的影响\cite{yang2023}，柔性模块接触模型则将冲击、结构变形与机器人运动联系起来\cite{wucontact2026}。这些研究说明，面向组装控制的模型应同时关注安装端运动传递、接触约束建立和载荷变化。

在模型应用层面，研究已由响应预测延伸到运动与控制设计。振动反馈利用机器人与结构的相互作用调节响应\cite{swei2020}；基于线性分式表示的路径优化将柔性动态和参数不确定性纳入性能评价\cite{rodrigues2025}；直接轨迹优化与递推动力学方法进一步结合耦合方程和作业约束改善求解效率\cite{liu2026,gong2026}。但高保真仿真、轨迹规划与高频反馈对模型信息的要求不同，计算效率的提高还需结合测量误差评价其闭环效果。面向实时控制，仍需提取安装界面运动与受力对末端响应的主要影响，明确简化模型的误差及适用范围。

\topic{（2）柔性支撑下机械臂的运动与阻抗控制}

柔性支撑条件下，基座运动通过机械臂运动链传递至作业端，使末端位姿保持与轨迹跟踪受到扰动。围绕机械臂运动控制，反作用零空间及其优化方法通过调整冗余关节运动减少安装端反作用\cite{nenchev1999,yin2023}；基于柔性状态估计与振动反馈的方法将支撑运动信息纳入控制，改善耦合条件下的末端响应\cite{patel2024,zhou2025}。这类方法为基座运动影响下的机械臂控制提供了途径，但控制效果受构型、剩余自由度、状态测量及关节驱动能力制约，需结合安装端运动与末端任务进行设计。

机械臂执行接触任务时，还需协调末端位姿跟踪与接触受力，通过阻抗设计建立运动偏差与力和力矩之间的动态关系。基于$\SE$的几何阻抗通过刚体位姿空间中的势能与误差描述协调平移和转动\cite{seo2023}；李群与李代数势能比较及阻抗映射研究，揭示了误差表示对有限位姿偏差下恢复作用的影响\cite{seo2024,kim2025}。针对模型与环境不确定性，滑模补偿和模糊增益调节用于减弱动力学失配对机械臂控制的影响\cite{liudb2024}，环境刚度估计与分层控制进一步结合柔性关节、基座振动及接触变化，调节机械臂末端的运动与柔顺响应\cite{liuadaptive2026}。

上述研究为柔性支撑下机械臂的运动与接触控制提供了基础，但末端参考运动和设计阻抗仍需通过有限的关节驱动实现。基座运动改变末端的实际位姿、速度与接触受力，关节运动或力矩约束激活又可能导致跟踪偏差，使实际柔顺响应偏离设计值。因此，需要围绕机械臂末端的运动精度与接触柔顺性，将安装端运动与受力、几何阻抗及关节力矩分配纳入统一分析，研究基座扰动和驱动约束共同作用下的位姿跟踪性能与阻抗可实现性。

\topic{（3）接触信息引导的位姿估计与柔顺组装}

柔顺组装正由利用接触反馈调节受力，向利用接触约束修正目标位姿发展。已有空间模块对接研究通过接触力误差调整柔顺轨迹，并完成地面验证\cite{shi2024}。但较小的接触力仍可能对应错误的接口偏置，目标位姿不确定时，需要同时利用受力反馈与配合几何。

接触几何利用主要包括联合状态估计与接触模型配准。前者融合触觉、机器人运动和接触约束，借助主动运动更新物体位姿与接触状态\cite{kimtactile2023}；后者通过探测采样局部接触流形，与预建接口模型配准获得位姿修正\cite{negi2025}。柔顺运动与主动接触交互的结合，使机器人能够逐步定位孔位和修正插接方向\cite{chen2025}。这些方法通过多次观测累积几何信息，其估计能力取决于接口形状、接触分布与探测运动。

在柔性支撑条件下，接触位姿变化同时包含目标偏差、基座振动及柔顺响应，用于估计的探测运动和参考更新又会改变接触载荷。因此，需分析基座测量与时序误差向目标配准的传播，协调参考修正与实际跟踪能力，明确估计精度改善能够转化为对准精度和接触载荷改善的条件。

综上，现有研究已形成耦合动力学、机械臂运动与阻抗控制及接触估计的方法基础。本项目进一步贯通安装界面表征、受约束阻抗实现与接触引导，研究各环节误差对组装响应的共同影响，并通过可控基座激励下的地面实验检验其适用范围。

\subsubsection{与相关指南的关系}

本项目主要对应指南2.2“轻质柔性结构动力学设计与分析”中柔性空间结构在轨组装及刚柔耦合体系力学控制研究，围绕机器人安装界面的动态作用、受约束运动与接触组装开展研究；同时衔接指南2.5关于多源不确定性条件下闭环分析与控制的研究内容。
\begingroup
\small
\setstretch{1.08}
\setlength{\parskip}{0pt}

\renewcommand{\refname}{主要参考文献}
\begin{thebibliography}{99}
\setlength{\itemsep}{0.15em}
\bibitem{hu2025} 胡海岩, 田强, 文浩, 罗凯, 马小飞. 极大空间结构在轨组装的动力学与控制[J]. 力学进展, 2025, 55(1): 1--29. DOI: \doi{10.6052/1000-0992-24-044}.
\bibitem{wu2024} 吴志刚, 蒋建平, 邬树楠, 李庆军, 王兴, 谭述君, 邓子辰. 航天结构空间组装动力学与控制研究进展[J]. 力学进展, 2024, 54(2): 344--390. DOI: \doi{10.6052/1000-0992-23-041}.
\bibitem{papadopoulos2021} Papadopoulos E, Aghili F, Ma O, Lampariello R. Robotic manipulation and capture in space: A survey[J]. Frontiers in Robotics and AI, 2021, 8: 686723. DOI: \doi{10.3389/frobt.2021.686723}.
\bibitem{sonneville2021} Sonneville V, Scapolan M, Shan M, Bauchau O A. Modal reduction procedures for flexible multibody dynamics[J]. Multibody System Dynamics, 2021, 51(4): 377--418. DOI: \doi{10.1007/s11044-020-09770-w}.
\bibitem{min2026} Min K, Fan W, Ren H, Zhou P, Wang Y, Zhang X. Dynamic modeling and model order-reduction of on-orbit assembled multi-plate structures[J]. International Journal of Mechanical Sciences, 2026, 313: 111296. DOI: \doi{10.1016/j.ijmecsci.2026.111296}.
\bibitem{ni2024} Ni S, Chen W, Chen T. Dynamic modeling and analysis of multi-flexible-link space manipulators under time-varying dynamic boundary conditions[J]. Advances in Space Research, 2024, 74(10): 5224--5243. DOI: \doi{10.1016/j.asr.2024.07.067}.
\bibitem{yang2023} Yang G, Zhang L, Yu S, Meng S, Wang Q, Li Q. Influences of space perturbations on robotic assembly process of ultra-large structures[J]. Nonlinear Dynamics, 2023, 111(11): 10025--10048. DOI: \doi{10.1007/s11071-023-08395-w}.
\bibitem{wucontact2026} Wu L, Wang M, Luo J, Zhang D. Dynamic modeling and analysis of space robots with an irrotational contact model for on-orbit assembly[J]. Aerospace Science and Technology, 2026, 176: 112117. DOI: \doi{10.1016/j.ast.2026.112117}.
\bibitem{swei2020} Swei S S M, Jenett B, Cramer N B, Cheung K. Modeling and control of robot-structure coupling during in-space structure assembly[C]//AIAA Scitech 2020 Forum. 2020: AIAA 2020-1544. DOI: \doi{10.2514/6.2020-1544}.
\bibitem{rodrigues2025} Rodrigues R, Preda V, Sanfedino F, Alazard D. Dynamics modeling and path optimization for the on-orbit assembly of large flexible structures using a multi-arm robot[J/OL]. CEAS Space Journal, 2025. DOI: \doi{10.1007/s12567-025-00675-y}.
\bibitem{liu2026} Liu P, Wang M, Ma W, Luo J, Liu C. Direct trajectory optimization of coupled constrained system formed by the space manipulator and the flexible structure[J]. Aerospace Science and Technology, 2026, 168: 111038. DOI: \doi{10.1016/j.ast.2025.111038}.
\bibitem{gong2026} Gong Y, Wen H. Recursive modeling and trajectory optimization of space manipulator mounted on flexible structure[J]. Applied Mathematical Modelling, 2026, 155: 116741. DOI: \doi{10.1016/j.apm.2026.116741}.
\bibitem{nenchev1999} Nenchev D N, Yoshida K, Vichitkulsawat P, Uchiyama M. Reaction null-space control of flexible structure mounted manipulator systems[J]. IEEE Transactions on Robotics and Automation, 1999, 15(6): 1011--1023. DOI: \doi{10.1109/70.817666}.
\bibitem{yin2023} 尹旺, 王翔, 王为, 刘冬雨. 柔性基座冗余机械臂的最优反作用控制[J]. 中国空间科学技术, 2023, 43(2): 144--154. DOI: \doi{10.16708/j.cnki.1000-758X.2023.0029}.
\bibitem{patel2024} Patel D, Damaren C J. Tip and vibration control of space robots using estimated flexible coordinates[J]. Frontiers in Space Technologies, 2024, 5: 1447545. DOI: \doi{10.3389/frspt.2024.1447545}.
\bibitem{zhou2025} Zhou W, Ye Z, Wu S, Li Y. Coupling dynamic modeling and vibration control of quadruped-robot and large space structure during on-orbit assembly[J]. Space Solar Power and Wireless Transmission, 2025, 2(1): 1--9. DOI: \doi{10.1016/j.sspwt.2025.02.002}.
\bibitem{seo2023} Seo J, Prakash N P S, Rose A, Choi J, Horowitz R. Geometric impedance control on $\SE$ for robotic manipulators[J]. IFAC-PapersOnLine, 2023, 56(2): 276--283. DOI: \doi{10.1016/j.ifacol.2023.10.1581}.
\bibitem{seo2024} Seo J, Prakash N P S, Choi J, Horowitz R. A comparison between Lie group- and Lie algebra-based potential functions for geometric impedance control[C]//2024 American Control Conference (ACC). 2024: 1335--1342. DOI: \doi{10.23919/ACC60939.2024.10644201}.
\bibitem{kim2025} Kim J, Sung M, Choi Y, Park J, Chung W K. Impedance control design framework using commutative map between $\SE$ and $\mathfrak{se}(3)$[J]. IEEE Transactions on Robotics, 2025, 41: 6193--6212. DOI: \doi{10.1109/TRO.2025.3619066}.
\bibitem{liudb2024} 刘东博, 陈力. 面向在轨组装服务的空间机器人动力学分析及抑振阻抗控制设计[J]. 力学季刊, 2024, 45(3): 688--696. DOI: \doi{10.15959/j.cnki.0254-0053.2024.03.008}.
\bibitem{liuadaptive2026} Liu P, Ma W, Wang M, Luo J, Liu C. Adaptive impedance control of a flexible-joint manipulator for space assembly operations under base vibrations and environmental uncertainty[J]. Advances in Space Research, 2026, 77(7): 8057--8074. DOI: \doi{10.1016/j.asr.2026.02.005}.
\bibitem{shi2024} 史玲玲, 肖晓龙, 张晓峰, 范立佳, 单明贺, 田强. 空间机器人在轨组装多模块单元的对接力控制与地面实验[J]. 力学学报, 2024, 56(3): 800--816. DOI: \doi{10.6052/0459-1879-23-458}.
\bibitem{kimtactile2023} Kim S, Jha D K, Romeres D, Patre P, Rodriguez A. Simultaneous tactile estimation and control of extrinsic contact[C]//2023 IEEE International Conference on Robotics and Automation (ICRA). 2023: 12563--12569. DOI: \doi{10.1109/ICRA48891.2023.10161158}.
\bibitem{negi2025} Negi A, Manyar O M, Penmetsa D K, Gupta S K. Accurate pose estimation using contact manifold sampling for safe peg-in-hole insertion of complex geometries[C]//2025 IEEE 21st International Conference on Automation Science and Engineering (CASE). 2025: 617--624. DOI: \doi{10.1109/CASE58245.2025.11163996}.
\bibitem{chen2025} Chen Y, Kimble K, Qian H H, Chanrungmaneekul P, Seney R, Hang K. Robust peg-in-hole assembly under uncertainties via compliant and interactive contact-rich manipulation[C]//Robotics: Science and Systems XXI. 2025. DOI: \doi{10.15607/RSS.2025.XXI.060}.
\bibitem{qi2025} 祁伊凡, 单明贺, 田强. 空间多柔体系统缩比等效动力学建模方法研究[J]. 中国科学: 物理学 力学 天文学, 2025, 55(2): 224518. DOI: \doi{10.1360/SSPMA-2024-0302}.
\bibitem{su2026} Su W, Shi L, Müller A, Shan M. Cooperative manipulation control with task-prioritized real-time optimization for free-floating dual-arm space robots[J]. Aerospace Science and Technology, 2026, 171: 111584. DOI: \doi{10.1016/j.ast.2025.111584}.
\end{thebibliography}
\endgroup

\subsection{项目的研究内容、研究目标及拟解决的关键科学问题}

\subsubsection{研究内容}

本项目面向柔性桁架支撑下机器人的精确运动与模块组装，采用理论分析、数值仿真与地面实验相结合的方法，围绕安装界面动力学传递、受约束阻抗响应及接触估计与执行协调，开展以下三项研究。

\Needspace{8\baselineskip}
\topic{（1）柔性桁架支撑下机械臂的耦合动力学建模与运动控制}

针对桁架振动与机械臂运动的双向耦合，以及关节驱动约束引起的阻抗实现偏差，从安装界面的运动与受力出发，研究耦合作用的反向动力学表征与几何阻抗优化控制，揭示运动精度、接触柔顺性与安装端反作用之间的关联。主要研究内容如下：
\begin{enumerate}[label=\arabic*)]
 \item 研究安装界面运动与受力的功共轭表征，建立以作业端为浮动基的反向动力学模型，揭示基座振动与关节运动对末端响应的耦合传递规律。
 \item 研究$\SE$几何阻抗与动力学优化的协同设计，将关节运动及力矩约束纳入统一求解，形成兼顾位姿跟踪与柔顺响应的受约束运动控制方法。
 \item 研究界面载荷建模偏差与约束激活下的阻抗可实现性，分析局部闭环稳定与误差有界条件，阐明运动精度、柔顺响应及安装端反作用的关联。
\end{enumerate}

\Needspace{8\baselineskip}
\topic{（2）目标位姿不确定条件下的接触对准与柔顺组装}

针对基座运动引起的接触观测偏差与参考执行失配，以接触流形的几何约束和接触反馈为依据，研究目标位姿的局部可辨识性与有界参考修正，揭示接触估计、参考更新和实际阻抗共同作用下的对准与受力规律。主要研究内容如下：
\begin{enumerate}[label=\arabic*)]
 \item 研究接触流形约束下六维目标位姿的局部可辨识性，分析基座测量与时序误差向接触场配准的传播规律，建立弱约束方向的先验融合方法。
 \item 研究五轴约束导纳与轴向接近运动的协调机制，保持参考端口耗散性与零输入保持性质，构建修正速率和累计偏移受限的柔顺组装策略。
 \item 研究目标估计、参考修正与受约束执行的耦合误差传递，阐明估计精度转化为对准精度和接触载荷改善的条件，形成引导与执行协调方法。
\end{enumerate}

\Needspace{8\baselineskip}
\topic{（3）柔性支撑条件下机器人运动与组装的地面实验验证}

针对振动环境下运动精度与接触响应的实验验证需求，依托基座振动模拟与机器人作业实验平台，研究安装界面运动复现与多源同步观测方法，通过分项对照和组装集成实验，检验耦合动力学描述、受约束控制及接触引导方法的有效性。主要研究内容如下：
\begin{enumerate}[label=\arabic*)]
 \item 研究安装界面动态边界的运动复现与载荷表征，构建多源同步观测与独立精度评价方法，明确基座激励对桁架振动影响的可表征范围。
 \item 开展位姿保持与轨迹跟踪对照实验，辨析界面载荷建模偏差与约束激活对运动精度和实际阻抗的影响，验证动力学模型与控制的实时性。
 \item 开展典型接口对准、插接与到位保持集成实验，分离接触配准和参考修正的性能贡献，评价振动环境下组装精度、接触载荷与方法适用范围。
\end{enumerate}

\Needspace{5\baselineskip}\subsubsection{研究目标}

面向大型空间结构在轨组装需求，建立基于安装界面运动与受力表征的反向动力学模型，揭示桁架振动与机械臂运动的耦合传递规律及驱动约束下阻抗响应的形成机制，形成几何阻抗与动力学优化相结合的运动控制方法；阐明接触几何约束下目标位姿的局部可辨识性及估计与执行误差对组装响应的共同影响，形成接触场配准、有界参考修正与受约束运动控制相协调的柔顺组装方法；建立安装端运动复现与独立测量评价的地面实验方法，完成运动控制和典型接口组装验证，明确所提方法在不同基座振动、末端负载及初始位姿偏差条件下的适用范围。

\subsubsection{拟解决的关键科学问题}

\topic{（1）机器人与柔性桁架耦合作用下受约束运动与阻抗响应的形成机制}

柔性桁架振动通过安装端运动与受力改变机械臂的动态响应，机械臂驱动又通过安装端反作用影响柔性桁架的振动响应，使作业端运动与安装端受力相互关联。面向实时控制，需要保留影响末端响应的主要耦合信息，而安装端作用力与力矩的近似描述会带来动力学偏差；关节运动范围与驱动力矩限制还可能使期望运动无法完全实现，导致实际阻抗响应偏离设计值。因此，安装端耦合作用如何经受约束力矩分配影响末端运动与柔顺响应，是本项目需要解决的关键科学问题。


\topic{（2）基座运动与接触约束共同作用下对准精度与柔顺组装响应的形成机制}

基座运动通过观测坐标与测量时序影响接触位姿估计，接触几何与观测分布又决定不同位姿方向的约束能力，使目标估计与参考修正相互关联。面向目标位姿不确定的组装过程，需要利用接触信息持续修正对准目标，而参考更新受到关节驱动约束与实际阻抗响应的共同限制；估计偏差和跟踪滞后还可能转化为额外接触载荷，使目标估计精度的改善难以直接对应于组装性能提升。因此，接触估计与参考修正如何经受约束运动执行影响对准精度与接触柔顺性，是本项目需要解决的关键科学问题。

\subsection{拟采取的研究方案及可行性分析}

\subsubsection{总体研究思路与技术路线}

对应三项研究内容，按照“界面耦合表征与运动控制、接触引导与受约束执行协调、地面实验验证”的技术路线展开。首先，以安装端运动与载荷建立反向动力学描述，结合几何阻抗与动力学优化研究驱动约束下的末端响应；其次，利用接触流形的局部几何信息估计目标位姿，分析观测误差传播并通过有界参考修正协调对准与接触受力；最后，复现安装端运动，开展模型校核、分项对照与组装集成实验，以实测结果检验误差分析和参数选取依据。技术路线如图\ref{fig:route}所示。

\begin{figure}[htbp]
\centering
\begin{tikzpicture}[
 block/.style={draw=blue!45!black,fill=blue!3,rounded corners=3pt,
   text width=12.4cm,align=center,inner sep=8pt,font=\small},
 flow/.style={-{Stealth[length=2.5mm]},thick,draw=blue!50!black},
 node distance=0.55cm]
 \node[block,fill=blue!8] (need) {\textbf{柔性桁架上的精确运动与柔顺组装}\\[3pt]
 基座振动\qquad 关节驱动约束\qquad 目标位姿偏差};
 \node[block,below=of need] (control) {\textbf{研究内容一\quad 耦合动力学与运动控制}\\[3pt]
 反向动力学\quad 几何阻抗\quad 动力学优化};
 \node[block,below=of control,draw=green!40!black,fill=green!3] (contact) {\textbf{研究内容二\quad 接触对准与柔顺组装}\\[3pt]
 接触流形表征\quad 接触场配准\quad 有界参考修正};
 \node[block,below=of contact,draw=violet!60!black,fill=violet!3] (test) {\textbf{研究内容三\quad 地面实验验证}\\[3pt]
 安装端运动复现与同步观测\\[2pt]
 位姿保持\quad 轨迹跟踪\quad 典型接口组装};
 \draw[flow] (need) -- (control);
 \draw[flow] (control) -- (contact);
 \draw[flow] (contact) -- (test);
\end{tikzpicture}
\caption{项目技术路线与验证关系}
\label{fig:route}
\end{figure}
\FloatBarrier

\subsubsection{反向动力学建模与运动控制}

以作业端为浮动基、实际安装端为虚拟末端，沿作业端向安装端重新组织运动链，建立反向动力学模型。通过安装界面的运动与受力计入桁架对机械臂的作用，分析基座振动、关节运动与末端负载的耦合传递。模型基本形式为：

{\color{black}\begin{equation}
 \bm M\dot{\bm\nu}+\bm h
 =\bm S^{\mathsf T}\bm\tau_j
 +\bm J_b^{\mathsf T}\bm F_b
 +\bm J_g^{\mathsf T}\bm F_e .
 \label{eq:reverse_dynamics}
\end{equation}}

其中，$\bm\nu$由作业端体速度和关节速度组成，$\bm M$与$\bm h$分别为惯性矩阵和包含科氏、离心及重力作用的偏置项；$\bm S$为关节选择矩阵，$\bm\tau_j$为关节驱动力矩；$\bm F_b$、$\bm F_e$及其雅可比矩阵$\bm J_b$、$\bm J_g$分别对应实际作业端和安装端。自由运动时$\bm F_b=0$，接触时由力传感测量计入；力旋量与速度采用功共轭坐标表达。

以现有安装端静力平衡近似为基线，结合实测基座运动与界面载荷，研究动态偏差随激励和负载的变化规律，形成控制所需的载荷描述与误差范围。地面实验中，模拟机械臂末端与作业机械臂基座固定连接，对应模型的实际安装端；通过指令、实测运动与受力对照，校核地面激励和控制模型的对应关系。

\begingroup\tolerance=9999\emergencystretch=8em
在$\SE$上统一描述实际位姿$T_b$与期望位姿$T_d$，采用局部误差坐标\allowbreak$\color{black}\bm\lambda=\operatorname{Log}(T_b^{-1}T_d)^\vee$，将几何阻抗映射为误差动态\cite{kim2025}：
\par\endgroup

\begin{equation}
 \bm A_\lambda\ddot{\bm\lambda}
 +\bm D_\lambda\dot{\bm\lambda}
 +\bm K_\lambda\bm\lambda
 =\bm\gamma .
 \label{eq:impedance}
\end{equation}

其中，$\bm A_\lambda$、$\bm D_\lambda$、$\bm K_\lambda$为经几何映射得到的等效惯量、阻尼和刚度矩阵，$\bm\gamma$为与误差坐标及相对运动方向一致的广义接触输入。结合轨迹前馈和速度映射生成作业端参考加速度，校核误差坐标与其时间导数的一致性。

将参考加速度嵌入动力学二次规划（QP），以参考跟踪、关节虚拟阻尼和力矩正则为目标，以反向动力学为等式约束，统一计入关节运动范围、力矩及其变化率限制。通过参考松弛处理期望运动不可实现的情形，研究约束激活、优化权重与实际阻抗的关系。以几何阻抗误差能量为基础，在优化可行、测量误差有界及参考足够平滑的条件下，分析典型约束激活区间内的局部稳定与误差界，结合受控扰动实验确定阻抗参数和参考变化速率的选取依据。

\subsubsection{基于接触流形表征与参考修正的柔顺组装}

以已有物理方向辅助接触场配准方法（PCF）为基础，以离线接触位姿约束场值，以功共轭变换后的物理方向约束局部斜率，结合全局场与局部拟合表征接口接触流形。在线以目标先验为初值，在统一参考系中求解六维位姿：
\begin{equation}
 \widehat T_\star=\underset{T\in\SE}{\arg\min}
 \sum_i\rho\!\left(F(T^{-1}T_i)\right).
 \label{eq:registration}
\end{equation}
其中，$T_i$为接触位姿，$F$为局部位姿坐标中的接触高度残差，$\rho$为鲁棒损失。结合实测基座位姿和时序修正观测坐标，按接口特征尺度统一平移与转动的局部度量，利用残差雅可比的奇异值及对应方向分析局部可辨识性，研究弱约束方向的先验融合与执行限幅，并通过扰动回放检验观测误差传播。

利用约束导纳捕获方法（RAC）调节两个横向与三个姿态通道，轴向接近单独处理。在正定导纳阻尼和包含零的闭凸增量约束下，依据共轭接触输入求解修正增量，限制变化速率与累计偏移，并校核有限转动步长，保持零输入保持性质及参考端口离散广义功关系。根据估计敏感程度设置导纳阻尼和修正边界，将位姿及其速度、加速度参考一致地传递至下层控制。

通过局部线性化与受控扰动，分析目标估计、参考修正及实际跟踪误差对到位精度和接触载荷的贡献，考察加速度松弛对误差传递的影响，形成参考修正与驱动能力的匹配方法。局部分析和参考子系统功关系用于解释作用机制，整机性能通过配准、参考修正及联合控制的对照实验评价。

\subsubsection{地面实验系统与分项、集成验证}

依托现有平台，由基座运动模拟机械臂驱动作业机械臂基座产生可重复的平移与转动，复现桁架振动引起的安装端运动。按平台工作范围和预实验结果设置激励幅值、频率与方向，以末端配重调整负载，以同一安装状态下的静止基座作为基线。同步采集基座与末端位姿、关节状态及两端力和力矩，分析激励传递的幅值和相位偏差，建立界面运动与载荷表征方法。传感器标定、外参测定和时序同步作为测量保障，末端精度采用独立于控制反馈的测量评价。

运动控制验证采用位姿保持、圆轨迹跟踪和接触扰动实验，结合仿真与实测对照，分析载荷描述及约束激活对实际阻抗的影响。保持任务参考和激励一致，核查实测基座运动，记录位置均方根误差、旋转测地误差、关节力矩及其变化率、安装端受力，并比较设计阻抗与力和运动实测响应；统计计算耗时分布、最大值和超周期比例，验证实时执行能力。每组正式对照至少重复5次，参数对照集中在代表性工况。

组装验证选取一种典型接口和代表性载荷，将目标固定于独立刚性基准。先在基座静止时设置基础柔顺控制、仅配准、仅参考修正和二者结合的四组对照，再在代表性平移或转动激励下重复关键对照。统一初始偏差、执行时间和保持窗口，记录原始估计、限幅参考和实际到位结果，分离估计与执行误差，评价到位精度、达标次数、接触力峰值及安装端反作用。

集成验证贯通接近、对准、插接与到位保持。正式对照前依据接口容差和测量精度固定工况及判据，成功与失败实验均纳入评价，调参数据不计入独立验证，报告绝对误差及重复实验离散程度。给定基座激励用于检验末端对振动扰动的适应能力；安装端反作用用于分析机器人对支撑结构的载荷作用，其降低不直接等同于桁架振幅降低。

\subsubsection{可行性分析}

\topic{（1）研究思路与理论方法可行}

本项目围绕基座运动下机械臂的运动精度与接触柔顺性，以浮动基机器人动力学为建模基础，将几何阻抗、动力学优化与接触引导相结合，通过理论分析、数值仿真和地面实验逐步验证。团队已有的建模方法、控制算法与对接实验为各项研究提供了基础。

\textbf{在耦合动力学建模方面：}团队已在自由漂浮双臂空间机器人协同操作研究中，建立包含浮动基座、机械臂与抓持目标的闭环系统运动学和动力学模型，描述基座与机械臂的动力学耦合及双臂协同操作约束\cite{su2026}。本项目沿用浮动基多体系统建模框架，以作业端为浮动基重新组织运动链，通过实际安装端的运动与载荷表征外部作用。前期已实现反向机械臂模型及安装端载荷的静力近似描述，可在此基础上结合实测基座运动与载荷校核动态偏差，为面向控制的耦合建模提供直接基础。

\textbf{在机械臂运动与阻抗控制方面：}几何阻抗为末端位置与姿态的柔顺响应提供统一描述，动力学二次规划为运动目标与关节驱动能力的协调提供求解手段。团队已开展基于分层二次规划的任务优先级实时优化控制研究，积累了多任务协调与物理约束处理的方法基础\cite{su2026}；前期已进一步将几何阻抗与动力学优化结合，在关节运动和力矩约束下求取控制指令，并用于位姿保持和圆轨迹跟踪实验。已有控制算法与实验基础，可支撑本项目研究基座扰动和建模偏差对末端跟踪及柔顺响应的影响，并通过闭环误差分析与对照实验检验驱动约束下的阻抗可实现性。

\textbf{在接触引导组装方面：}团队已形成物理方向辅助接触场配准（PCF）与约束导纳捕获（RAC）方法，利用接触几何与物理方向信息估计目标位姿，通过有界参考修正实现柔顺对准，并完成仿真与多轮真机对接实验。现有方法具备修正速率与累计偏移约束、零输入保持性质，并已建立参考端口离散广义功关系，为接触引导过程的误差分析与柔顺控制提供了依据。本项目可在已有方法上引入基座运动补偿，分析观测偏差与参考修正的误差传递，再由静止基座逐步扩展到振动工况，检验接触引导与机械臂力矩控制的协调效果，具备明确的方法基础和实验验证途径。

\topic{（2）实验方案与研究条件可行}

\textbf{从实验方案看：}模拟机械臂与作业机械臂的固定连接可形成可控、可重复的安装端运动环境，配重与接口夹具可设置负载和初始偏差。通过基座与末端位姿、关节状态及两端受力的同步观测，可分别评价模型误差、跟踪性能与接触响应；结合独立刚性基准和精度测量，能够对照检验各控制环节，覆盖主要研究问题。

\textbf{从研究基础与条件看：}团队已有基座振动模拟与机器人作业平台、开放力矩控制程序、运动捕捉系统、IMU和六维力测量条件，并积累了多模块对接实验经验。拟完善的连接组件、测点、时序同步和接口夹具可依托现有设备、共享条件与项目预算落实。按单项模型与控制校核、接口对照和振动环境集成验证逐步推进，工作量与两年研究周期相匹配。

\subsection{本项目的特色与创新之处}

本项目以安装界面的运动与受力贯通耦合动力学、几何阻抗和接触组装，结合可控基座激励与独立测量，研究柔性支撑条件下运动精度与接触柔顺性的协调关系。特色与创新主要体现在以下两个方面。

\topic{（1）安装界面耦合作用与驱动约束共同作用下的阻抗实现方法}

针对柔性支撑条件下设计阻抗与实际响应不一致的问题，在反向动力学框架中将安装界面运动与载荷引入作业端响应分析，揭示界面载荷建模偏差与驱动约束激活对阻抗特性的共同影响；将几何阻抗目标与受约束力矩分配相结合，研究阻抗可实现性及局部误差界，形成依据界面扰动和驱动能力协调运动精度与柔顺响应的控制方法。

\topic{（2）基座运动下接触几何估计与受约束执行相协调的柔顺组装方法}

针对目标估计精度提升难以直接转化为组装性能改善的问题，将接触流形的局部可辨识性与实际阻抗响应联系起来，揭示观测误差、参考修正和执行偏差向对准误差及接触载荷的传递规律；结合弱约束方向先验、有界导纳修正及下层跟踪能力，形成接触引导与受约束执行相协调的方法，提高典型接口在基座振动条件下的对准精度与接触柔顺性。

\subsection{年度研究计划及预期研究成果}

\subsubsection{年度研究计划}

\topic{2027年1月至6月}
围绕研究内容一，建立安装界面耦合作用的反向动力学描述，校核几何阻抗与动力学坐标映射；同步开展研究内容三的安装端运动复现与多源观测，确定激励和负载工况，完善连接组件、测点及执行时序。形成动力学模型、界面基线数据和实验评价方案。

\topic{2027年7月至12月}
完成几何阻抗与动力学优化控制设计，分析建模偏差及约束激活下的阻抗可实现性，开展位姿保持、轨迹跟踪与阻抗响应实验；启动接触观测误差分析和动态接触场适用性检验。形成论文初稿1篇及年度报告，与重点实验室开展专题交流。

\topic{2028年1月至6月}
完成接触位姿局部可辨识性与误差传播分析，研究弱约束方向先验及参考修正与驱动能力的协调，将PCF与RAC接入运动控制；完成静止基座四组对照及代表性单向激励实验，形成方法论文与发明专利申请材料。

\topic{2028年7月至12月}
完成接触引导与受约束执行协调方法验证，开展运动与组装集成实验，评价不同基座激励和初始偏差下的组装精度、接触载荷及方法适用范围；整理程序、数据和测试报告，完成论文发表或录用、专利申请、成果交流与结题。

\subsubsection{预期研究成果与学术交流}

对应三项研究内容，形成安装界面耦合动力学表征与受约束阻抗响应的分析结果，开发包含反向动力学、几何阻抗与动力学优化的机器人运动控制程序1套，以及典型接口的接触配准与有界参考修正模块1套；建立安装端运动复现、同步观测与独立精度评价的地面实验流程，形成同步测量数据集，完成振动环境下运动与典型接口组装的地面原理验证，提交技术总结报告和实验测试报告各1份。围绕新增研究发表或录用学术论文3篇，申请发明专利1项，培养参与研究的研究生3名。

每年至少与重点实验室研究人员开展1次实质性学术交流，结合模型、实验和阶段成果组织研讨与学术报告；拟参加1次相关国内学术会议，并利用线上形式开展国际学术讨论。项目成果署名、资助标注和知识产权安排按重点实验室开放课题管理要求执行。

\section{研究基础与工作条件}

\subsection{工作基础}

团队围绕空间在轨操作，已形成柔性多体动力学、受约束机器人控制和接触对接的研究积累，可支撑本项目三项研究内容的衔接与实施。

\topic{（1）柔性多体动力学与模型简化基础}

申请人参与了柔性多体系统模态缩减研究\cite{sonneville2021}，团队进一步开展空间多柔体系统缩比等效建模\cite{qi2025}与几何精确梁动力学降阶研究，积累了柔性结构动力学分析和模型校核经验。上述研究积累可支撑本项目对机器人与柔性桁架相互作用的分析、安装端作用力与力矩的校核，以及地面试验条件与实际作业工况之间的对应分析。

\topic{（2）反向动力学建模与运动控制基础}

团队针对七自由度作业机械臂，已实现以作业端为浮动基、通过实际安装端作用力与力矩描述外部作用的反向动力学模型，并建立了几何阻抗与动力学优化相结合的运动控制程序。该程序利用关节状态以及运动捕捉系统与惯性测量单元（IMU）的测量数据生成运动参考，在关节运动、力矩及其变化约束下求取控制指令，已支持位姿保持、圆轨迹跟踪及固定基座轨迹跟踪试验，并具备数据记录与回放功能。现有安装端受力模型采用静力平衡近似，为本项目在模拟装置施加基座激励的条件下校核安装端动态作用力与力矩、研究末端运动与安装端力和力矩的关系提供了可扩展的实现基础。

\topic{（3）接触几何与柔顺对接基础}

团队已开展在轨组装柔顺控制与轨迹规划研究，并完成多模块对接力控制地面实验\cite{shi2024}。近期研究已形成PCF六维目标配准与RAC约束导纳捕获方法，建立了有界修正、零输入保持及参考端口离散广义功关系。论文主表选取的五轮真机试验中，完整方法在65次试验中58次满足两秒保持窗口内位置与姿态均方根误差不超过0.5~mm和0.5$^\circ$的判据，仅RAC为35/65；同组60段非零偏差配准回放的平均三维位置与姿态误差为2.085~mm和1.622$^\circ$。现有评分采用示教对齐基准，真机通过导纳与位置伺服执行。上述结果为本项目进一步开展动态基座下的力矩控制融合及独立测量验证提供了基础。

\subsection{工作条件}

\topic{（1）已具备的条件}

本项目依托现有基座振动模拟与机器人作业实验平台，通过基座运动模拟装置向作业机械臂安装端施加振动扰动。现有作业机械臂的开放力矩控制程序、运动捕捉系统与IMU的状态输入接口及同步记录接口，可支持运动控制集成与对照实验；已有模块对接平台、六维力/力矩测量及接触记录，可支撑接口标定和柔顺接触研究。作业机械臂的运动控制程序及接触记录回放流程为运动控制实现与接触对准方法验证提供了基础。

依托单位另建有约$4\,\mathrm m\times10\,\mathrm m$的在轨操作微重力模拟地面实验平台，可为相关地面验证提供配套条件。本项目悬挂装置的实验结果按照地面重力下的基座振动工况进行分析。

\topic{（2）尚需完善的条件及解决途径}

围绕基座振动下的闭环验证，拟完善模拟装置与作业机械臂之间的悬挂连接法兰、转接组件及可换配重，布置基座运动与安装端力和力矩测点，完成位姿测量系统、力和力矩传感器及传感器空间外参的标定，并实现模拟装置与作业机械臂的执行时序同步。结合现有对接平台完善典型接口夹具与独立位姿测量，使基座激励、末端运动和最终对接精度能够分别评价。连接组件、接口改装与专项测试按项目预算落实，计算与通用测量设备优先共享。

拟依托重点实验室在航空航天结构力学及控制方面的研究条件，围绕安装端作用力与力矩校核、运动与受力测量及控制实验开展联合讨论和测试，按阶段研究进度落实设备共享与试验安排。


\subsection{申请人简介}

\topic{（1）学历与研究工作简历}

单明贺，北京理工大学教授、博士生导师，国家级青年人才，主要从事空间在轨操作动力学与控制、柔性多体动力学和空间机器人技术研究。2007年9月至2011年7月在哈尔滨工业大学学习，获飞行器制造工程学士学位；2011年9月至2013年7月在哈尔滨工业大学学习，获航空宇航制造工程硕士学位；2013年11月至2018年6月在荷兰代尔夫特理工大学学习，获航天工程博士学位。2018年11月至2020年11月在美国马里兰大学帕克分校从事博士后研究，2021年1月起在北京理工大学工作。

申请人主持国家科技重大专项子课题、国家自然科学基金面上项目和青年项目等科研项目10余项，发表高水平论文30余篇，获授权发明专利7项。相关研究覆盖柔性结构动力学、空间机器人操作与地面验证，为本项目的耦合建模、控制方法研究和实验组织提供支撑。

\Needspace{7\baselineskip}
\topic{（2）近期与本项目有关的主要论著}

\begin{enumerate}[label=\arabic*.]
 \item \begin{minipage}[t]{\linewidth}祁伊凡, 单明贺, 田强. 空间多柔体系统缩比等效动力学建模方法研究[J]. 中国科学: 物理学 力学 天文学, 2025, 55(2): 224518. DOI: \doi{10.1360/SSPMA-2024-0302}.
 \end{minipage}
 \item \begin{minipage}[t]{\linewidth}史玲玲, 肖晓龙, 张晓峰, 范立佳, 单明贺, 田强. 空间机器人在轨组装多模块单元的对接力控制与地面实验[J]. 力学学报, 2024, 56(3): 800--816. DOI: \doi{10.6052/0459-1879-23-458}.
 \end{minipage}
 \item \begin{minipage}[t]{\linewidth}Yifan Qi, Shilei Han, Minghe Shan, Mingwu Li, Qiang Tian. Invariant-manifold-based model reduction for geometrically exact beam dynamics[J]. Computer Methods in Applied Mechanics and Engineering, 2026, 450: 118643. DOI: \doi{10.1016/j.cma.2025.118643}.
 \end{minipage}
\end{enumerate}

% 接触研究稿为匿名研究材料，作者名单未提供，已在“工作基础”说明。
% 不将匿名稿件冒列为申请人已发表论著，正式论著清单按实际发表状态更新。
\topic{（3）学术奖励}

\begin{enumerate}[label=\arabic*.]
\item 单明贺，国际空间研究委员会青年科学家最佳论文奖（COSPAR Outstanding Paper Award for Young Scientists），2022年。对应论文为：Minghe Shan, Jian Guo, Eberhard Gill. An analysis of the flexibility modeling of a net for space debris removal[J]. Advances in Space Research, 2020, 65(3): 1083--1094. DOI: \doi{10.1016/j.asr.2019.10.041}。
\item 张大羽，罗凯，赵将，韩石磊，朱佳龙，单明贺，王辉，刘凡，薛永刚，王立朋. 空间大型柔性可展结构高可靠展开与超低频模态辨识关键技术及应用. 中国振动工程学会科学技术奖二等奖，2023年度。
\end{enumerate}

\subsection{项目组主要成员简介}

项目组由单明贺、苏文康、黄朕、李融冰和任启明组成，围绕耦合建模、控制方法与地面验证拟作如下分工。

\textbf{单明贺：}项目负责人，研究基础及代表性成果见前节。负责总体研究方案与进度组织，统筹机器人与柔性桁架的耦合建模、闭环性能分析及与重点实验室的合作交流。

\textbf{苏文康：}已开展自由漂浮空间机器人任务优先协同操作与实时优化控制研究。拟承担几何阻抗与动力学优化控制设计、驱动约束协调及基座振动下的末端运动与安装端力和力矩分析。

\textbf{黄朕：}已开展大型空间望远镜多机器人组装序列规划研究，具备面向在轨组装任务组织操作流程的研究基础。拟承担典型作业过程与参考轨迹设计、实验工况组织及结果评价，服务于本项目单机器人运动与对接验证。

\textbf{李融冰：}拟承担作业机械臂的控制程序集成与实物实验，重点开展实时力矩接口实现、基座振动下运动控制验证和实验数据分析。

\textbf{任启明：}拟承担接触配准与参考修正方法衔接、典型接口柔顺组装实验及误差分析。

\Needspace{7\baselineskip}
\topic{近期与本项目有关的主要论著}

\begingroup
\small
\setstretch{1.12}
\begin{enumerate}[label=\arabic*.]
\item \begin{minipage}[t]{\linewidth}Wenkang Su, Lingling Shi, Andreas Müller, Minghe Shan. Cooperative manipulation control with task-prioritized real-time optimization for free-floating dual-arm space robots[J]. Aerospace Science and Technology, 2026, 171: 111584. DOI: \doi{10.1016/j.ast.2025.111584}.
\end{minipage}
\item \begin{minipage}[t]{\linewidth}Zhen Huang, Lingling Shi, Lijia Fan, Yan Yang, Minghe Shan. Multirobot Assembly Sequence Planning of a Large Space Telescope[J]. Journal of Spacecraft and Rockets, 2025, 62(3): 915--927. DOI: \doi{10.2514/1.A36129}.
\end{minipage}
\item \begin{minipage}[t]{\linewidth}Qiming Ren, Minghe Shan. A unified framework for compliant control and trajectory planning in robotic in-orbit assembly[J]. Acta Astronautica, 2026, 243: 32--45. DOI: \doi{10.1016/j.actaastro.2026.01.029}.
\end{minipage}
\end{enumerate}
\endgroup


\subsection{近五年承担科研项目情况}

本项目依托既有研究基础，重点研究柔性安装界面的动态作用、受约束阻抗响应及接触参考修正，形成相应模型、算法和典型接口实验结果。

\needinfo{近五年科研项目的正式名称、编号、经费来源、起止年月及承担角色，并按实际任务书说明与本项目的联系和任务区别。}

\clearpage
\section{经费预算}

% 重点课题20--30万元，本稿按30万元编制；数量和单价为需求测算，非供应商报价。
申请经费30.00万元，围绕模拟装置与作业机械臂的悬挂连接、接口改装、同步测量与控制接口、专项测试及成果产出安排。依托现有地面实验平台及计算设备开展研究，优先共享可用仪器设备。经费预算如下。

\begingroup
\small
\begin{longtable}{P{3.15cm}>{\centering\arraybackslash}p{1.35cm}P{9.6cm}}
\toprule
预算科目 & 金额（万元） & 计算根据及理由 \\
\midrule
\endfirsthead
\toprule
预算科目 & 金额（万元） & 计算根据及理由 \\
\midrule
\endhead
\midrule
\multicolumn{3}{r}{续下页}\\
\endfoot
\bottomrule
\endlastfoot
\textbf{合计} & \textbf{30.00} & 服务于耦合动力学与运动控制、接触组装和地面实验验证三项研究内容。 \\
设备费 & 4.00 & 同步采集和实时控制接口改装所需设备及配件1套，按4.00万元测算，用于模拟装置与作业机械臂的状态、基座与末端运动及力和力矩测量数据的时序同步。 \\
材料费 & 6.00 & 悬挂连接法兰、转接组件及配套材料按3.60万元测算；典型接口夹具及可换配重材料1.40万元；传感安装、线缆等耗材1.00万元。 \\
测试化验加工费 & 9.00 & 悬挂连接及接口夹具加工装配3批，每批1.50万元，共4.50万元；基座与末端位姿测量、安装端力和力矩传感器标定3批，每批0.80万元，共2.40万元；外协动态测量与联合测试3批，每批0.70万元，共2.10万元。共享仪器免费使用部分不重复计费。 \\
差旅费 & 3.00 & 合作试验、专题交流及参加国内学术会议等约10人次，按每人次0.30万元测算。 \\
会议费 & 0.50 & 项目专题研讨及实验方案评议2次，按每次0.25万元测算。 \\
国际合作与交流费 & 0.00 & 以线上学术交流为主，不安排出国经费。 \\
出版/文献/信息传播/知识产权事务费 & 3.50 & 论文出版相关费用2篇，每篇按1.00万元测算，共2.00万元；发明专利申请及相关事务1.00万元；文献和信息资料0.50万元。 \\
专家咨询费 & 1.00 & 动力学建模及实验方案咨询约5人次，按每人次0.20万元测算；占总经费3.33\%。 \\
劳务费 & 3.00 & 参与建模、控制设计和实验工作的研究生2人，每人累计15个月，每人每月0.10万元；占总经费10.00\%。 \\
\end{longtable}
\endgroup

预算按研究需求测算，设备与测试项目结合现有条件核定，相关支出按实际发生及适用标准执行。专家咨询费和劳务费分别满足申请书模板规定的不超过总经费5\%和10\%的比例要求。

% ==================== 报告正文结束 ====================
\end{document}
